\subsection{Punto de Slater y cota de los multiplicadores del tope}\label{sec:P01}
Del Lema~\ref{lem:maximizador} sólo se usa la existencia y unicidad del maximizador canónico;
no se presupone su interioridad ni se afirma que coincida con las probabilidades de ocupación de las
casillas en una solución elegida uniformemente al azar.
Recordamos que $\mathcal D$ es la unión de las dos familias diagonales del tablero,
con $|\mathcal D|=4n-2$, y que
\[
 m_d(p)=\sum_{c\in d}p_c,\qquad h_d(p)=1-m_d(p).
\]
El politopo $K_n$ y el funcional $H_n$ son los de~\eqref{eq:Kn} y~\eqref{eq:Hn}, con la convención
$0\log0=0$. Denotamos por $p=\pi^{(n)}$ el único maximizador de $H_n$ en $K_n$.

\paragraph{Un punto de Slater para $n\ge32$.}
Sean $r_i=(i-1/2)/n$ y
\[
 q_{ij}=\frac1n\left[1-\frac15\cos(2\pi r_i)\cos(2\pi r_j)\right].
\]
La suma de los cosenos de los $n$ puntos medios es cero para $n\ge2$;
por tanto las sumas de filas y columnas de $q$ son exactamente uno y
$4/(5n)\le q_{ij}\le6/(5n)$.
La densidad estrictamente admisible $\nu^\sharp(r,c)=1-\frac15\cos(2\pi r)\cos(2\pi c)$
de~\cite[Observación~2.4]{Perfil} tiene, para $0\le d\le1$, la marginal
\[
 m_D[\nu^\sharp](d)=1-d-\frac15\left[
       \frac{1-d}{2}\cos(2\pi d)-\frac{\sin(2\pi d)}{4\pi}\right].
\]
Su derivada es a lo sumo $-1+(1+\pi)/5<0$, luego $m_D[\nu^\sharp]\le9/10$.
Las simetrías dan la misma cota en las dos familias completas.
Cada suma diagonal de $q$ es una cuadratura de puntos medios de
$\nu^\sharp$ en la fibra correspondiente, con extremos múltiplos de $1/n$.
La longitud paramétrica de la fibra es a lo sumo uno y su constante
de Lipschitz es a lo sumo $4\pi/5$. El error de cuadratura es, pues,
a lo sumo $(4\pi/5)/(2n)\le1/20$ para $n\ge32$, usando $\pi<4$.
En consecuencia, con $M_0=6/5$ y $\kappa=1/20$,
\[
 q_{ij}\le M_0/n<1/8,\qquad h_d(q)\ge\kappa.       \tag{P1.1}
\]
Se obtiene así un punto de Slater en el espacio afín de las sumas de filas y columnas.

\paragraph{Positividad en presencia del tope.}
Tomamos $p(t)=(1-t)p+tq$. Si $p_c=0$, su entropía espacial aporta
$-q_c t\log t+O(t)$ a $H_n(p(t))-H_n(p)$. Si $h_d(p)=0$, su
entropía diagonal aporta $-h_d(q)t\log t+O(t)$. Las demás entropías
aportan $O(t)$; también lo hace $3(1-p)\log(1-p)$, cuya derivada
está acotada en $[0,1/4]$. Si se da alguna de esas igualdades, la
diferencia es $-At\log t+O(t)>0$ para $t>0$ pequeño, con $A>0$,
lo que contradice la maximalidad. Por ello
\[
 p_{ij}>0,\qquad h_d(p)>0.                         \tag{P1.2}
\]

\paragraph{Condiciones de optimalidad con multiplicadores del tope.}
Por (P1.2), el funcional es diferenciable cerca de $p$ y las únicas desigualdades
potencialmente activas son $p_{ij}\le1/4$.
La condición de Slater proporciona multiplicadores finitos $\alpha_i,\beta_j$ y
$\mu_{ij}\ge0$. Una de las igualdades de filas y columnas es redundante y se suprime
al aplicar las condiciones de Karush--Kuhn--Tucker (KKT)~\cite[cap.~5]{BV}. Si $D(i,j),A(i,j)$ son las
dos diagonales de la casilla, resulta
\[
 -\log p_{ij}-3\log(1-p_{ij})+\log h_{D(i,j)}+\log h_{A(i,j)}-2
 =\alpha_i+\beta_j+\mu_{ij},\quad
 \mu_{ij}(1/4-p_{ij})=0.                          \tag{P1.3}
\]
Sean $x_{ij}=np_{ij}$, $\ell_d=-\log h_d(p)\ge0$ y
$d_0=3\log(4/3)$. Multiplicando (P1.3) por $q_{ij}-p_{ij}$ y
sumando se cancelan las igualdades y la constante $-2$. Se obtiene
\[
 0=E+F+Q+U,                                      \tag{P1.4}
\]
\[
 E=\sum_{ij}(p_{ij}-q_{ij})\log x_{ij},\qquad
 F=3\sum_{ij}(p_{ij}-q_{ij})\log(1-p_{ij}),
\]
\[
 Q=\sum_{d\in\mathcal D}(h_d(q)-h_d(p))\ell_d,\qquad
 U=\sum_{ij}\mu_{ij}(1/4-q_{ij}).
\]
En particular $U$ lleva signo positivo: la derivada hacia $q-p$
es $-U$, por complementariedad.

\paragraph{Cota global que incluye los multiplicadores del tope.}
La convexidad da $\sum p_{ij}\log(np_{ij})\ge0$. Además
\[
 \sum q_{ij}\log x_{ij}\le\frac{M_0}{n}\sum\log^+x_{ij}
 \le\frac{M_0}{n}\sum x_{ij}=M_0n,
\]
de modo que $E\ge-M_0n$. Como $\log(1-p)\le0$,
\[
 F\ge3\sum p_{ij}\log(1-p_{ij})\ge-d_0n.
\]
La desigualdad $h(-\log h)\le1/e$ y (P1.1) dan
\[
 Q\ge\kappa\sum_d\ell_d-4n/e,\qquad
 U\ge\frac18\sum_{ij}\mu_{ij}.
\]
Por (P1.4), si $C_0=M_0+d_0+4/e$,
\[
 \kappa\sum_d\ell_d+\frac18\sum_{ij}\mu_{ij}\le C_0n.
\]
Como $\kappa\le1/8$, con $T=C_0/\kappa$ queda
\[
 \frac1n\left(\sum_d\ell_d+\sum_{ij}\mu_{ij}\right)\le T.    \tag{P1.5}
\]
Esta estimación no presupone que ningún multiplicador $\mu$ sea cero.
